El ensayo se hace con el filtro cargado únicamente por la punta del
osciloscopio (\SI{10}{\mega\ohm} en paralelo con \SI{12}{\pico\farad}), que a
estas frecuencias no carga nada. La Figura~\ref{fig:banco} muestra la cadena.

\begin{figure}[H]
\centering
\begin{tikzpicture}[
    font=\small,
    caja/.style={draw=black!45, rounded corners=3pt, minimum height=1.5cm,
                 minimum width=2.9cm, align=center, fill=black!3},
    dut/.style={caja, draw=xtalPrimary, line width=1pt, fill=xtalPrimary!6},
    flecha/.style={-{Stealth[length=5pt]}, thick},
  ]
  \node[caja]              (gen)  {Generador\\\footnotesize \SI{1}{\volt} p.\,a p.};
  \node[dut,  right=1.5cm of gen] (dut)  {Filtro RLC\\\footnotesize bajo ensayo};
  \node[caja, right=1.5cm of dut] (osc)  {Osciloscopio\\\footnotesize 2 canales};
  \node[caja, below=0.9cm of dut] (arf)  {Analizador de respuesta\\\footnotesize en frecuencia};

  \draw[flecha, xtalInput]  (gen) -- node[above, font=\footnotesize] {$v_i$} (dut);
  \draw[flecha, xtalOutput] (dut) -- node[above, font=\footnotesize] {$v_o$} (osc);
  \draw[flecha, xtalInput]  (arf.west) -| ([xshift=-0.75cm]dut.west);
  \draw[flecha, xtalOutput] ([xshift=0.75cm]dut.east) |- (arf.east);
  \node[font=\footnotesize\itshape, below=0.1cm of arf] {barrido de \SI{20}{\hertz} a \SI{100}{\kilo\hertz}};
\end{tikzpicture}
\caption{Cadena de medición. El mismo circuito se ensaya de dos maneras: barrido
de frecuencia con el analizador, y escalón con el generador y el osciloscopio.}
\label{fig:banco}
\end{figure}

\begin{table}[H]
\centering
\caption{Instrumental y condiciones del ensayo.}
\label{tab:instrumental}
\small
\begin{tabular}{lll}
\toprule
\textbf{Ensayo} & \textbf{Instrumento} & \textbf{Condiciones} \\
\midrule
Respuesta en frecuencia & Analizador de barrido & 10 puntos/década, \SI{20}{\hertz}--\SI{100}{\kilo\hertz} \\
Respuesta al escalón    & Generador + osciloscopio & \SI{100}{\hertz}, 500 muestras, \SI{500}{\micro\second} por división \\
Valor de componentes    & Puente LCR & \SI{1}{\kilo\hertz}, cuatro hilos \\
\bottomrule
\end{tabular}
\end{table}
